% !TeX root = ../main.tex
%% ------------------------------------------------------------------------- %%
\chapter{Introdução}
\label{cap:introducao}

%% ------------------------------------------------------------------------- %%
\section{Contextualização e Descrição do Problema}
\label{sec:contextualizacao}

O mercado global de Controladores Lógicos Programáveis (CLPs) tem apresentado crescimento expressivo nos últimos anos, impulsionado pelos investimentos crescentes em automação industrial e pela modernização de sistemas de controle legados: estima-se que o setor movimenta atualmente cerca de US\$ 13,33 bilhões, com projeção de expansão contínua até 2031 \cite{mordor:2026}.
Esse crescimento reflete a consolidação da automação industrial como pilar central da Indústria 4.0, na qual controladores programáveis assumem papel crítico na interconexão de máquinas, sensores e processos produtivos.
Nesse cenário, contudo, a expansão do setor convive com uma forte fragmentação tecnológica: a predominância de soluções proprietárias e de arquiteturas de \textit{hardware} específicas por fabricante ainda limita a portabilidade de algoritmos de controle e favorece o aprisionamento tecnológico dos usuários finais.
É nesse contexto que se insere a proposta deste trabalho: uma máquina virtual (\textit{VM}) embarcada, capaz de interpretar instruções de lógica \textit{Ladder} --- baseada na norma IEC 61131-3 \cite{iec:2025} --- de maneira independente da plataforma de \textit{hardware} subjacente.

A lógica \textit{Ladder} (\textit{Ladder Diagram}) é uma das cinco linguagens de programação padronizadas pela IEC 61131-3 para controladores industriais, caracterizada pela representação gráfica de contatos e bobinas que remete aos antigos diagramas elétricos a relés, o que a tornou historicamente a linguagem mais difundida na automação industrial.
Uma máquina virtual, por sua vez, corresponde a uma camada de \textit{software} capaz de interpretar um conjunto de instruções (\textit{bytecode}) de forma independente do \textit{hardware} subjacente --- de modo análogo ao papel desempenhado por máquinas virtuais de propósito geral em outros domínios da computação.
Reunir esses dois conceitos exige a definição de uma Arquitetura de Conjunto de Instruções (\textit{ISA} --- \textit{Instruction Set Architecture}) customizada e de tamanho fixo (32 \textit{bits}), projetada para otimizar o processo de busca e decodificação (\textit{fetch-decode}) e garantir o determinismo do tempo de varredura; a pesquisa foca especificamente na migração dessa arquitetura para microcontroladores da família \textit{STM32} (núcleo \textit{ARM Cortex-M7}), explorando técnicas de abstração de \textit{hardware} (\textit{HAL}) e metodologias de \textit{Model-Based Design}.
Propõe-se ainda a segregação estrutural entre o motor de execução (\textit{firmware} estático) e a lógica de aplicação (\textit{bytecode} tratado como dado), estabelecendo um ambiente de execução isolado (\textit{sandbox}) que previne falhas críticas de processamento e viabiliza a gestão segura do ciclo de vida da aplicação industrial.

Apesar do crescimento do setor, a rigidez das arquiteturas convencionais de CLPs impõe barreiras significativas à inovação, dificultando a portabilidade de algoritmos de controle para plataformas de \textit{hardware} mais modernas e acessíveis sem a necessidade de reengenharia completa do \textit{firmware}.
Além disso, a validação de lógicas de controle complexas diretamente no \textit{hardware} final apresenta riscos operacionais e custos elevados, carecendo de metodologias que integrem eficientemente a simulação computacional com a execução física.

No âmbito das soluções de código aberto atualmente disponíveis, como o ecossistema \textit{OpenPLC} \cite{alves:2018}, a abordagem predominante baseia-se na compilação de diagramas \textit{Ladder} para código C nativo, seguida de compilação estática.
Embora eficiente em termos de tempo de execução, essa metodologia funde a lógica de controle aos \textit{drivers} de \textit{hardware} em um binário monolítico, de modo que a adaptação para novas plataformas ou pinagens exige modificações diretas no código-fonte dos \textit{drivers} em baixo nível, limitando drasticamente a portabilidade.
Adicionalmente, a execução do código traduzido diretamente no processador elimina barreiras de proteção inerentes a sistemas operacionais: um erro na modelagem da lógica do usuário pode resultar em falhas críticas do sistema, comprometendo a estabilidade do equipamento industrial.

Por outro lado, abordagens alternativas baseadas em interpretadores de linguagens de \textit{script} genéricas, como o \textit{MicroPython} \cite{micropython:2025}, falham em atender aos requisitos de tempo real estrito da automação industrial, devido à gestão não determinística de memória (\textit{Garbage Collection}), que introduz variações inaceitáveis no tempo de varredura.
Diante do exposto, o problema técnico enfrentado por este trabalho consiste na ausência de uma camada de virtualização eficiente e aberta que permita a execução segura e determinística de lógicas de controle padronizadas em microcontroladores de propósito geral, garantindo a fidelidade comportamental entre o modelo projetado em ambiente de simulação e o sistema embarcado final.


%% ------------------------------------------------------------------------- %%
\section{Objetivos}
\label{sec:objetivo}

\subsection{Objetivo Geral}
\label{sec:objgeral}

Diante do problema apresentado, este trabalho tem como objetivo geral desenvolver e validar uma máquina virtual embarcada capaz de interpretar \textit{bytecode} de lógica \textit{Ladder} com fidelidade semântica à norma IEC 61131-3, executando de forma determinística e agnóstica de \textit{hardware} em um microcontrolador da família \textit{STM32}.

\subsection{Objetivos Específicos}
\label{sec:objespecificos}

\begin{itemize}
	\item Especificar e documentar formalmente uma Arquitetura de Conjunto de Instruções (\textit{ISA}) de 32 \textit{bits}, compatível com as primitivas operacionais do subconjunto \textit{Ladder Diagram} da IEC 61131-3;
	\item Modelar lógicas de controle de referência em ambiente \textit{Simulink}/\textit{Stateflow}, segundo a metodologia de \textit{Model-Based Design};
	\item Implementar, em linguagem C, o motor de execução (\textit{VM}) responsável pelo ciclo de busca, decodificação e execução do \textit{bytecode};
	\item Desenvolver uma Camada de Abstração de \textit{Hardware} (\textit{HAL}) para o microcontrolador \textit{STM32F7} (núcleo \textit{ARM Cortex-M7});
	\item Validar cruzadamente o comportamento da \textit{VM} embarcada frente ao modelo de referência simulado, comparando as tabelas de imagem de saída para um mesmo conjunto de estímulos de entrada;
	\item Analisar quantitativamente o desempenho da solução, mensurando o tempo de varredura (\textit{scan cycle}) e o consumo de memória \textit{RAM} e \textit{Flash}.
\end{itemize}


%% ------------------------------------------------------------------------- %%
\section{Justificativa}
\label{sec:justificativa}

A relevância deste trabalho decorre da lacuna identificada entre as soluções de código aberto disponíveis e os requisitos de portabilidade e determinismo da automação industrial.
Ao segregar o motor de execução da lógica de aplicação, a arquitetura proposta elimina a necessidade de recompilação do \textit{firmware} a cada alteração da lógica \textit{Ladder}, reduzindo o tempo de comissionamento e o risco de falhas decorrentes de atualizações completas de memória.
O caráter de ambiente isolado (\textit{sandbox}) confere ainda uma camada de segurança inexistente nas soluções que traduzem a lógica do usuário diretamente para código nativo, contendo eventuais erros estruturais antes que comprometam a estabilidade do sistema embarcado.

Do ponto de vista social e econômico, a proposta de uma solução aberta e portável para múltiplas plataformas de \textit{hardware} contribui para reduzir o aprisionamento tecnológico a fabricantes específicos, democratizando o acesso a tecnologias de automação hoje concentradas em soluções proprietárias de alto custo.


%% ------------------------------------------------------------------------- %%
\section{Estrutura do Documento}
\label{sec:organizacao}

A metodologia adotada caracteriza-se como pesquisa aplicada de natureza experimental e abordagem predominantemente quantitativa, estruturada em etapas sequenciais com entregas incrementais que partem da especificação formal da arquitetura até a validação funcional integrada ao \textit{hardware}; o detalhamento de cada etapa --- especificação da \textit{ISA}, modelagem de referência, implementação do motor de execução, desenvolvimento da \textit{HAL}, integração com \textit{hardware}, validação cruzada e análise de desempenho --- será apresentado nos capítulos de desenvolvimento deste trabalho.

O restante deste trabalho está organizado da seguinte forma: o Capítulo~\ref{cap:fundamentacao} apresenta a fundamentação teórica, contemplando a norma IEC 61131-3, os trabalhos relacionados e o estado da prática em soluções para execução de lógicas de controle industrial, além dos conceitos de arquitetura que sustentam a solução proposta.
Os capítulos subsequentes apresentarão a concepção e a modelagem detalhada da solução, os resultados obtidos na validação cruzada e na análise de desempenho, seguidos das considerações finais e das sugestões para trabalhos futuros, no Capítulo~\ref{cap:finais}.
